\begingroup
\renewcommand{\arraystretch}{1.45}
\Needspace{45\baselineskip}
\InterpolationArticle{28}
\textit{Example.} In the scholarly correspondence edited by the distinguished Baron \textsc{de Zach}, Volume X, p. 188, there is a table giving the northern and southern limits of the zodiac of \textit{Pallas}. The declination of each limit is regarded as a periodic function of right ascension, which advances in that table in steps of five degrees. To illustrate the preceding investigations, we shall apply them here to the northern limit; its declination will therefore be our $X$ and its right ascension our $x$. From that table we extract the following period of 12 terms:
{\small\setlength{\arraycolsep}{3pt}\[
\begin{array}{r|r r l c r}
x&\multicolumn{5}{c}{X}\\ \hline
0^{\mathord{0}}&6^{\mathord{0}}&48'&\text{N.}&=+&408'\\
30&1&29&.\,.\,.\,.\,.\,.\,.&+&89\\
60&1&6&\text{S.}\,.\,.&-&66\\
90&0&10&\text{N.}\,.\,.\,.&+&10\\
120&5&38&.\,.\,.\,.\,.\,.\,.&+&338\\
150&13&27&.\,.\,.\,.\,.\,.\,.&+&807\\
180&20&38&.\,.\,.\,.\,.\,.\,.&+&1238\\
210&25&11&.\,.\,.\,.\,.\,.\,.&+&1511\\
240&26&23&.\,.\,.\,.\,.\,.\,.&+&1583\\
270&24&22&.\,.\,.\,.\,.\,.\,.&+&1462\\
300&19&43&.\,.\,.\,.\,.\,.\,.&+&1183\\
330&13&24&.\,.\,.\,.\,.\,.\,.&+&804
\end{array}
\]}
\SourceLeaf{318}
\Needspace{47\baselineskip}
Let us first divide this period into three periods of four terms each:
{\small\setlength{\arraycolsep}{3pt}\[
\begin{array}{r r|r r|r r}
a=0^{\mathord{0}}&A=+408&a'=30^{\mathord{0}}&A'=+89&a''=60^{\mathord{0}}&A''=-66\\
b=90^{\mathord{0}}&B=+10&b'=120^{\mathord{0}}&B'=+338&b''=150^{\mathord{0}}&B''=+807\\
c=180^{\mathord{0}}&C=+1238&c'=210^{\mathord{0}}&C'=+1511&c''=240^{\mathord{0}}&C''=+1583\\
d=270^{\mathord{0}}&D=+1462&d'=300^{\mathord{0}}&D'=+1183&d''=330^{\mathord{0}}&D''=+804
\end{array}
\]}
In the formula
\[
\begin{aligned}
&X'={}\gamma+\gamma'\cos x+\gamma''\cos2x\\[5pt]
&+\delta'\sin x+\delta''\sin2x
\end{aligned}
\]
the following values therefore result:
\[
\begin{aligned}
&\gamma=\tfrac14(A+B+C+D)=779'{,}5\\[5pt]
&\gamma'=\tfrac12(A\cos a+B\cos b+C\cos c+D\cos d)\\[5pt]
&=\tfrac12(A-C)=-415'{,}0\\[5pt]
&\delta'=\tfrac12(A\sin a+B\sin b+C\sin c+D\sin d)\\[5pt]
&=\tfrac12(B-D)=-726'{,}0\\[5pt]
&\gamma''=\tfrac14(A\cos2a+B\cos2b+C\cos2c+D\cos2d)\\[5pt]
&=\tfrac14(A-B+C-D)\\[5pt]
&=+43'{,}5\\[5pt]
&\delta''=\tfrac14(A\sin2a+B\sin2b+C\sin2c+D\sin2d)=0
\end{aligned}
\]
and similarly for the second and third periods. This gives
{\small\setlength{\arraycolsep}{3pt}\[
\begin{array}{l|r|r|r|r|r|r}
\text{For the period}&\text{where }y=4x&\gamma&\gamma'&\delta'&\gamma''&\delta''\\ \hline
\text{first}&0^{\mathord{0}}&+779{,}5&-415{,}0&-726{,}0&+43{,}5&0\\
\text{second}&120^{\mathord{0}}&+780{,}2&-404{,}5&-721{,}4&+9{,}9&+17{,}1\\
\text{third}&240^{\mathord{0}}&+782{,}0&-413{,}5&-713{,}3&+11{,}7&-20{,}3
\end{array}
\]}
\Needspace{42\baselineskip}
Here, furthermore, $\gamma$ is represented by the formula
\[
\begin{aligned}
&\tfrac13(779'{,}5+780{,}2+782{,}0)\\[5pt]
&+\tfrac23(779{,}5+780{,}2\cos120^{\mathord{0}}+782{,}0\cos240^{\mathord{0}})\cos4x\\[5pt]
&+\tfrac23(780{,}2\sin120^{\mathord{0}}+782{,}0\sin240^{\mathord{0}})\sin4x
\end{aligned}
\]
or by
\[
\begin{aligned}
&780{,}6-1{,}1\cos4x-1{,}0\sin4x,\quad\text{and similarly}
\end{aligned}
\]
\[
\begin{aligned}
&\gamma'\quad\text{by}\quad-411{,}0-4{,}0\cos4x+5{,}2\sin4x\\[5pt]
&\delta'\quad\text{by}\quad-720{,}2-5{,}8\cos4x-4{,}7\sin4x\\[5pt]
&\gamma''\quad\text{by}\quad+21{,}7+21{,}8\cos4x-1{,}1\sin4x\\[5pt]
&\delta''\quad\text{by}\quad-1{,}1+1{,}1\cos4x+21{,}6\sin4x
\end{aligned}
\]
On substituting these values in $X'$, the following formula representing the 12 proposed values results:
\SourceLeaf{319}
\[
\begin{aligned}
&+780{,}6\\[5pt]
&-411{,}0\cos x-720{,}2\sin x\\[5pt]
&+43{,}4\cos2x-2{,}2\sin2x\\[5pt]
&-4{,}3\cos3x+5{,}5\sin3x\\[5pt]
&-1{,}1\cos4x-1.0\sin4x\\[5pt]
&+0{,}3\cos5x-0{,}3\sin5x\\[5pt]
&+0{,}1\cos6x
\end{aligned}
\]
\Needspace{42\baselineskip}
Let us next divide the same period into four periods of three terms each, each of which must therefore be satisfied by a formula of the first order,
\[
\begin{aligned}
&X'=\gamma+\gamma'\cos x+\delta'\sin x
\end{aligned}
\]
Here the following values are found:
\begin{center}
\begin{tabular}{r|c|r|r|r}
For the period & where $y=3x$ & $\gamma$ & $\gamma'$ & $\delta'$\\\hline
first & $0^{0}$ & $+776{,}3$ & $-368{,}3$ & $-718{,}8$\\
second & $90^{0}$ & $+786{,}0$ & $-414{,}5$ & $-676{,}0$\\
third & $180^{0}$ & $+785{,}0$ & $-453{,}0$ & $-721{,}1$\\
fourth & $270^{0}$ & $+775{,}0$ & $-408{,}2$ & $-765{,}0$
\end{tabular}
\end{center}
From this the following formulas representing $\gamma,\gamma',\delta'$ are derived:
\[
\begin{aligned}
&\gamma\quad.\,.\,.\,.\quad+780{,}6-4{,}3\cos3x+5{,}5\sin3x\\[5pt]
&\qquad{}+0{,}1\cos6x\\[5pt]
&\gamma'\quad.\,.\,.\,.\quad-411{,}0+42{,}3\cos3x-3{,}2\sin3x\\[5pt]
&\qquad{}+0{,}3\cos6x\\[5pt]
&\delta'\quad.\,.\,.\,.\quad-720{,}2+1{,}2\cos3x+44{,}5\sin3x\\[5pt]
&\qquad{}+0{,}3\cos6x
\end{aligned}
\]
On substituting these values in $\gamma+\gamma'\cos x+\delta'\sin x$ and treating the last parts in $\gamma,\gamma',\delta'$ that contain $\cos6x$ according to the rules of Article 26, precisely the same formula is obtained as the one we reached above.
\endgroup
